\documentclass[11pt]{article}
\usepackage[margin=0.85in]{geometry}
\usepackage{amsmath,amssymb,amsthm,mathtools,microtype}
\usepackage[dvipsnames]{xcolor}
\usepackage[colorlinks=true,linkcolor=MidnightBlue,citecolor=MidnightBlue,urlcolor=MidnightBlue]{hyperref}
\newtheorem{theorem}{Theorem}
\newtheorem{lemma}[theorem]{Lemma}
\newtheorem{remark}[theorem]{Remark}
\newcommand{\M}{\mathcal M}
\newcommand{\E}{\mathcal E}
\newcommand{\hpE}{\mathsf h^c_{p,E}}
\newcommand{\hpat}{\mathsf h^c_{p,\mathrm{at}}}
\newcommand{\hpc}{\mathsf h^c_p}
\newcommand{\tr}{\tau}
\newcommand{\supp}{\operatorname{supp}}
\title{A Boyd-index reduction for symmetric atomic Hardy spaces}
\author{Autonomous proof attempt for arXiv:2201.10219}
\date{11 August 2026}

\begin{document}
\maketitle

\begin{abstract}
Hong--Ma--Wang ask whether their symmetric atomic Hardy space
$\mathsf h^c_{p,E}(\M)$ equals the conditioned Hardy space $\mathsf h^c_p(\M)$
for $0<p<1$, and whether the intervening crude symmetric-atomic space also
coincides.  We prove that, whenever the lower Boyd index satisfies $p_E>2$,
both symmetric spaces equal the standard $(p,2)$-atomic space, with equivalent
quasi-norms.  Thus, throughout this large class of symmetric spaces, the
question is exactly equivalent to the foundational unresolved problem
$\mathsf h^c_{p,\mathrm{at}}=\mathsf h^c_p$.  The proof combines the source's
factorization identities with the constructive $(p,\infty)$-atomic theorem of
Chen--Randrianantoanina--Xu.
\end{abstract}

\section{Statement and proof intuition}

Let $(\M,\tr)$ be a finite von Neumann algebra with normalized trace and a
fixed filtration $(\M_n)$.  Let $E$ be a separable symmetric Banach function
space on $(0,\infty)$ with $1<p_E\le q_E<\infty$.  We use the definitions in
\cite{HMW}: a $(p,E)_c$-atom $a$ has $\E_n(a)=0$, right support below some
$e\in\M_n$, and
\begin{equation}\label{eq:Eatom}
 \|a\|_{\mathsf h_E^c}
 \le \tr(e)^{1-1/p}\|e\|_{E^\times}^{-1}.
\end{equation}
The spaces built from these atoms and from the crude factorizations in
\cite[Remark 4.13]{HMW} are denoted by $\mathsf h^c_{p,E}$ and
$\mathsf h^c_{p,E;\mathrm{crude}}$.  The usual $(p,2)$-atomic space is
$\hpat$.

\begin{theorem}[Boyd-index reduction]\label{thm:main}
Let $0<p\le1$ and $2<p_E\le q_E<\infty$.  Then
\[
 \mathsf h^c_{p,E}(\M)
 =\mathsf h^c_{p,E;\mathrm{crude}}(\M)
 =\mathsf h^c_{p,\mathrm{at}}(\M)
\]
with equivalent quasi-norms.  The row analogues follow by adjoints, and so do
the corresponding componentwise mixed-space identities.
\end{theorem}

\paragraph{Idea of the proof.}
There are two complementary factorizations.  Since $p_E>2$, the source's
Boyd-index lemma yields $L_2=E\odot F$ for a symmetric space $F$; on an
$e$-supported operator this converts the $E$-bound \eqref{eq:Eatom} exactly
into the $(p,2)$ scaling.  Conversely, choosing $q>q_E$ gives a factorization
$E=G\odot L_q$, so every $(p,q)$ atom is an $E$-atom.  The constructive theorem
of \cite{CRX} decomposes every $(p,2)$ atom into $(p,\infty)$ atoms, closing the
reverse inclusion.  The same $L_2$ factorization and a two-norm interpolation
of the adapted factor in a crude symmetric atom reduce it to a standard crude
$(p,2)$ atom.

\section{The two atom comparisons}

We use two standard consequences of the pointwise-product formulas quoted and
proved in Section 3 of \cite{HMW}.  Equivalent renormings only change the
constants below.

\begin{lemma}[Local factorization estimates]\label{lem:local}
Assume $p_E>2$ and write $t=\tr(e)$ for a projection $e$.
\begin{enumerate}
\item There is a symmetric $F$ such that $L_2=E\odot F$ and, for $z=ze$,
\begin{equation}\label{eq:Eto2}
 \|z\|_2\le C_E\|z\|_E\frac{t^{1/2}}{\|e\|_E}.
\end{equation}
\item For every $q>q_E$ there is a symmetric $G$ such that
$E=G\odot L_q$ and, for $z=ze$,
\begin{equation}\label{eq:qtoE}
 \|z\|_E\le C_{E,q}\|z\|_q\frac{\|e\|_E}{t^{1/q}}.
\end{equation}
\end{enumerate}
Moreover $\|e\|_E\|e\|_{E^\times}=t$.
\end{lemma}

\begin{proof}
Because $p_E>2$, Lemma 3.2 of \cite{HMW} (Lemma 2.2 in its source file)
renorms $E^{(1/2)}$ as a symmetric Banach space.  Apply Lozanovskii
factorization to $E^{(1/2)}$ and then square the product identity to obtain
$L_2=E\odot F$.  The fundamental-function product formula gives
$\|e\|_E\|e\|_F\asymp t^{1/2}$.  Factoring $z=ze$ gives
\eqref{eq:Eto2}.  For $q>q_E$, Lemma 3.3 of \cite{HMW} gives
$E=G\odot L_q$; now
$\|e\|_G t^{1/q}\asymp\|e\|_E$, and $z=ez$ gives
\eqref{eq:qtoE}.  Finally $L_1=E\odot E^\times$ gives the last identity.
\end{proof}

If $a$ is a $(p,E)_c$-atom supported by $e$, then its conditioned square
function is also supported by $e$ \cite[Lemma 4.7]{HMW}.  Lemma
\ref{lem:local}, \eqref{eq:Eatom}, and the last projection identity give
\begin{align*}
 \|a\|_{\mathsf h_2^c}
 &\le C_E\|a\|_{\mathsf h_E^c}\frac{t^{1/2}}{\|e\|_E}\\
 &\le C_E\frac{t^{1-1/p}}{\|e\|_{E^\times}}
              \frac{t^{1/2}}{\|e\|_E}
 =C_E t^{1/2-1/p}.
\end{align*}
Thus every $(p,E)$ atom is a uniformly bounded multiple of a $(p,2)$ atom,
and
\begin{equation}\label{eq:forward}
 \mathsf h^c_{p,E}\hookrightarrow\hpat.
\end{equation}

For the converse, fix $q>q_E$.  If $a$ is a $(p,q)_c$-atom supported by $e$,
then \eqref{eq:qtoE} yields
\[
 \|a\|_{\mathsf h_E^c}
 \le C_{E,q}\|e\|_E t^{-1/q}t^{1/q-1/p}
 =C_{E,q}\|e\|_E t^{-1/p}
 =C_{E,q}\frac{t^{1-1/p}}{\|e\|_{E^\times}}.
\]
So every $(p,q)$ atom is a bounded multiple of a $(p,E)$ atom.  By
\cite[Theorem 4.3]{CRX}, every $(p,2)$ atom is a sum of $(p,\infty)$ atoms with
uniformly controlled $\ell_p$ coefficients.  On a support of trace $t$, an
$\mathsf h_\infty^c$ bound implies the required $\mathsf h_q^c$ bound by
Hölder, so these are also $(p,q)$ atoms.  Consequently
\begin{equation}\label{eq:reverse}
 \hpat\hookrightarrow\mathsf h^c_{p,E}.
\end{equation}
Equations \eqref{eq:forward}--\eqref{eq:reverse} prove the first and third
spaces in Theorem \ref{thm:main} are equal.

\section{The crude symmetric atoms}

It remains to control the middle space.  For $p=1$, \cite{HMW} records
$\mathsf h^c_{1,E}=\mathsf h^c_1$ and
\[
 \mathsf h^c_{1,E}\subset
 \mathsf h^c_{1,E;\mathrm{crude}}\subset\mathsf h^c_1,
\]
so equality is immediate.  Assume $0<p<1$, and put
\[
 u=\frac p{1-p},\qquad s=\frac{2p}{2-p},
 \qquad \frac1p=\frac12+\frac1s.
\]
Let $a=yb$ be a crude $(p,E)_c$-atom.  Thus $\E_n(y)=0$,
$\|y\|_{\mathsf h_E^c}\le1$, and
\begin{equation}\label{eq:crude}
 b\in L_u(\M_n)\cap E^\times(\M_n),\qquad
 \|b\|_u\|b\|_{E^\times}\le1.
\end{equation}
Taking $e=1$ in \eqref{eq:Eto2} gives
$\|y\|_{\mathsf h_2^c}\le C_E$.  Choose $q>q_E$.  The continuous inclusion
$L_q\hookrightarrow E$ dualizes to $E^\times\hookrightarrow L_{q'}$, hence
$\|b\|_{q'}\le C_{E,q}\|b\|_{E^\times}$.  Define $v$ by
\[
 \frac2v=\frac1u+\frac1{q'}=\frac1p-\frac1q.
\]
Noncommutative Hölder applied to $|b|^2$ gives
\begin{equation}\label{eq:binterp}
 \|b\|_v^2=\||b|^2\|_{v/2}
 \le\|b\|_u\|b\|_{q'}
 \le C_{E,q}.
\end{equation}
But
\[
 \frac1s-\frac1v=\frac{1-p}{2p}+\frac1{2q}>0,
\]
so $v>s$.  Since $\tr(1)=1$, $L_v\hookrightarrow L_s$ contractively.
Equations \eqref{eq:crude}--\eqref{eq:binterp} show that $a$ is a uniformly
bounded multiple of a standard crude $(p,2)_c$-atom.  Proposition 3.9 of
\cite{CRX} decomposes every such crude atom into $(p,2)$ atoms with controlled
$\ell_p$ coefficients.  Therefore
\[
 \mathsf h^c_{p,E;\mathrm{crude}}\hookrightarrow\hpat.
\]
Together with the defining inclusion
$\mathsf h^c_{p,E}\subset\mathsf h^c_{p,E;\mathrm{crude}}$ and
\eqref{eq:reverse}, this completes the proof of Theorem \ref{thm:main}.

\section{What remains open}

For $0<p<1$, \cite{CRX} proves
$\mathsf h^c_{p,\mathrm{alg}}=\mathsf h^c_p$ but explicitly leaves open
$\hpat=\mathsf h^c_p$.  The weak decomposition there controls coefficients
only in $\ell_1$, whereas the atomic quasi-norm requires $\ell_p$ control.
Theorem \ref{thm:main} therefore gives, for every $E$ with $p_E>2$,
\[
 \mathsf h^c_{p,E}=\mathsf h^c_p
 \quad\Longleftrightarrow\quad
 \mathsf h^c_{p,\mathrm{at}}=\mathsf h^c_p.
\]
This is a substantial reduction but not a solution of the foundational
right-hand equality.  At the endpoint $p_E=2$, the proof still applies whenever
the structural factorization $L_2=E\odot F$ with matching fundamental
functions is available (in particular for $E=L_2$); the Boyd index alone no
longer supplies it.

\begin{thebibliography}{9}
\bibitem{HMW}
G. Hong, C. Ma, and Y. Wang,
\emph{John--Nirenberg inequalities for noncommutative column BMO and
Lipschitz martingales}, arXiv:2201.10219v1 (2022).

\bibitem{CRX}
Z. Chen, N. Randrianantoanina, and Q. Xu,
\emph{Atomic decompositions for noncommutative martingales},
arXiv:2001.08775v1 (2020).
\end{thebibliography}

\end{document}

